\section{Late Breaking}
\sectionkicker{Post-Cutoff / Friday 18:00 ET 以後}

\textbf{この欄は通常本文へ時系列を混ぜないためのbufferである。} W32のEditorial Cutoff後に重要度が上がった2件を、短い速報として分離する。詳細な再検証は次号へ持ち越す余地を残す。

\begin{latebreaking}[OpenAI Astra --- Critical cyber capabilityを「排除できない」]
8月7日、OpenAIはAstraの直近internal evaluationを受け、Preparedness Framework上の\textbf{Critical cyber capabilityを排除できない}と公表した\cite{openai-astra-critical-cyber}。これはAstraが正式にCriticalへ分類済みという意味ではない。OpenAIはassessment継続中としつつ、isolated testing、network/tool restriction、weight protection、monitoring、sandboxingを強化し、\textbf{強化後のcontrol requirementを満たさないAstra関連internal activityをpauseした}と説明している。

OpenAIが示すCritical thresholdは、多数のhardened real-world critical systemに対するautonomous zero-day exploit development、またはhigh-level goalからnovel end-to-end cyberattack strategyを遂行する能力を含む\cite{openai-astra-critical-cyber}。今号p.\pageref{sec:astra-lead}の数学Lead Storyとは同じAstraでも別eventであり、8月1日のresearch resultへ遡って混ぜない。
\end{latebreaking}

\begin{latebreaking}[SGLang v0.5.17 --- Kimi K3 / MiniMax H3をserving stackへ]
8月8日、SGLang v0.5.17が公開され、Kimi K3とMiniMax H3のday-0 supportが追加された\cite{sglang-0517}。H3ではT2VA / FL2VA / Ref2VAを含むvideo+audio生成path、Kimi K3ではreasoning/tool-call/OpenAI-compatible servingを含むmodel-specific supportが入る。

このreleaseの重要性はversion numberそのものより、今号で扱った2つのlarge/open modelが\textbf{発表から短期間でserving frameworkの具体的なexecution pathへ落ちた}ことにある。SGLangはH3についてB200、H100、2×RTX 5090のvalidated deployment profileも提示しており、model releaseとdeployment ecosystemの距離が縮まっている。
\end{latebreaking}

\begin{claimboundary}[Late BreakingはMain windowの事実へ逆流させない]
Astra cyber updateとSGLang v0.5.17はいずれもCutoff後である。今号のMain Storyを補足するが、「Cutoff時点ですでに判明していた」「Main-window community reactionだった」とは扱わない。
\end{claimboundary}
